\documentclass[a4paper,11pt]{article}
\usepackage{exam-style-842-unit}
\examinehead{842 信息化概论 · 第 1 单元练习册}{练习一\quad 概念与辨析}
\begin{document}

\unitcover{842 信息化概论 · 单元练习册}{第 1 单元\quad 信息技术与信息社会}{练习一\quad 概念与辨析（共 100 分）}{约 70 分钟}{数据、信息、知识、智慧；信息的分类与特点；信息科学三论；信息技术与信息产业；信息社会与信息化；四次信息技术革命；我国信息化发展战略。}

\begin{knowmap}
\textbf{认知链：}数据（记载符号、信息的载体）$\rightarrow$ 信息（加工后能减少不确定性的内容）$\rightarrow$ 知识（规律化、逻辑化的认识）$\rightarrow$ 智慧（从不同知识导出的一般性原理）\\[2pt]
\textbf{信息分类：}按加工程度 一次（原始）$\rightarrow$ 二次（文摘·索引）$\rightarrow$ 三次（综述·年鉴）；按表现形式 文献型·档案型·统计型·动态型·图像型；按来源 书本·报刊·人·具体事物\\[2pt]
\textbf{三大要素与信息特点：}物质 · 能源 · 信息；信息有普遍性·无限性·客观性，以及可共享、可存储、可传输、可扩散、可转换、可度量、可压缩等特点\\[2pt]
\textbf{信息科学：}信息论（香农 1948）· 控制论（维纳 1948）· 系统论（贝塔朗菲）$\Rightarrow$ 以电子计算机等为主要工具\\[2pt]
\textbf{信息技术四范畴：}传感（感觉器官）· 通信（神经系统）· 计算机（思维器官）· 控制（效应器官）$\Rightarrow$ 支撑起信息产业
\end{knowmap}

\begin{keyterms}
\textbf{数据}：记载客观事物性质、状态和相互关系的物理符号，是信息的载体。\\
\textbf{信息}：数据加工后能为某个目的使用的内容，可减少不确定性。\\
\textbf{知识}：使信息从定量到定性得以实现的、抽象的、具有逻辑的东西。\\
\textbf{智慧}：从完全不同的知识中导出的一般性原理。\\
\textbf{一次信息}：未经加工的原始信息，如口头、图片、原始记录。\\
\textbf{二次信息}：对一次信息加工后得到的规则有序信息，如文摘、索引。\\
\textbf{三次信息}：系统组织、压缩、分析一次和二次信息的结果，如综述、年鉴。\\
\textbf{信息科学}：研究信息现象及其运动规律和应用方法的科学。\\
\textbf{信息论}：研究信息的产生、获取、变换、传输、存储、处理识别及利用。\\
\textbf{控制论}：维纳 1948 年创立，方法有功能模拟法和黑箱法。\\
\textbf{系统论}：贝塔朗菲奠基，研究一切客观现实系统共同的特征与规律。\\
\textbf{信息技术}：研究信息的产生、获取、存储、传输、处理及应用的技术。\\
\textbf{信息产业}：进行信息的收集、整理、存储、传输、处理及其应用服务的产业。
\end{keyterms}

\begin{contrast}{数据 / 信息 / 知识 / 智慧（每行依次为：核心定义·产生方式·典型用途）}
\citem{数据}{客观事物性质、状态和相互关系的记载符号}{原始记录、尚未加工}{田间观测值、某年出生人数}
\citem{信息}{数据加工后能为某个目的使用的内容}{对数据加工、诠释}{婴儿潮的缘由、病虫害简报}
\citem{知识}{使信息从定量到定性得以实现的、抽象且具逻辑的东西}{对信息归纳、规律化}{病虫害发生规律、栽培模型}
\citem{智慧}{从完全不同的知识中导出的一般性原理}{跨领域知识综合与情境判断}{决定是否改种、何时灌溉}
\end{contrast}

\begin{contrast}{一次 / 二次 / 三次信息（按产生先后与加工深度划分；每行依次为：定义·典型形式·使用价值）}
\citem{一次信息}{未经加工的原始信息}{口头、图片、数字、表格、清单}{反映现场真实情况，但零散、难检索}
\citem{二次信息}{对一次信息加工处理后得到的信息}{文摘、索引、数据卡片}{规则有序，易于存储、检索、传递和使用}
\citem{三次信息}{系统组织、压缩和分析一次、二次信息的结果}{综述、专题报告、辞典、年鉴}{深入研究的结晶，可直接支撑决策}
\end{contrast}

\begin{contrast}{信息科学 / 信息技术 / 信息产业（每行依次为：定位·研究或业务对象·关键标志）}
\citem{信息科学}{研究信息现象及其运动规律和应用方法的科学}{以信息论、控制论、系统论为理论基础}{以电子计算机等为主要工具}
\citem{信息技术}{扩展人类信息器官功能的技术}{信息的产生、获取、存储、传输、处理及应用}{传感、通信、计算机、控制四大范畴}
\citem{信息产业}{进行信息收集、整理、存储、传输、处理及其应用服务的产业}{信息产品与信息服务的生产经营}{战略性、高渗透、知识智力密集等特征}
\end{contrast}

\begin{contrast}{信息论 / 控制论 / 系统论（每行依次为：创始人或代表作·研究内容·主要方法或观点）}
\citem{信息论}{香农，1948 年《通信的数学原理》等}{信息的产生、获取、变换、传输、存储、处理识别及利用}{狭义以编码理论为中心；广义即信息科学}
\citem{控制论}{维纳，1948 年出版《控制论》}{机器与生物中的控制和通信的一般规律}{功能模拟法、黑箱法}
\citem{系统论}{贝塔朗菲}{一切客观现实系统共同的特征、本质、原理和规律}{事物自成系统、互为系统，是动态稳定的开放系统}
\end{contrast}

\qtypename{一、单项选择题}{10 题，每题 3 分，共 30 分}
\q 在信息科学中，对客观事物的性质、状态和相互关系等进行记载的物理符号称为（\quad）
\choicesline{A. 信息}{B. 数据}{C. 知识}{D. 智慧}

\q 信息论的创始人香农认为，信息是（\quad）
\choices{A. 消除不确定性的东西}{B. 与外部世界交换的内容}{C. 物质系统运动的本质特征}{D. 从知识导出的一般性原理}

\q 使信息从定量到定性的过程得以实现的、抽象的、具有逻辑的东西是（\quad）
\choicesline{A. 数据}{B. 信息}{C. 知识}{D. 智慧}

\q 对一次信息进行加工处理后得到的、规则有序的信息是（\quad）
\choicesline{A. 一次信息}{B. 二次信息}{C. 三次信息}{D. 动态信息}

\q 综述、专题报告、辞典、年鉴等属于（\quad）
\choicesline{A. 一次信息}{B. 二次信息}{C. 三次信息}{D. 动态信息}

\q 按信息的（\quad）划分，信息可分为文献型、档案型、统计型、动态型和图像型。
\choicesline{A. 加工程度}{B. 表现形式}{C. 来源}{D. 时效长短}

\q 下列不属于信息基本特点的是（\quad）
\choicesline{A. 可共享性}{B. 可压缩性}{C. 可度量性}{D. 可消耗性}

\q 维系人类社会存在及发展的三大要素是（\quad）
\choices{A. 物质、能源、信息}{B. 材料、资金、技术}{C. 物质、能量、知识}{D. 信息、知识、智慧}

\q 计算机技术是信息的处理和存储技术，它对应于人的（\quad）
\choicesline{A. 感觉器官}{B. 神经系统}{C. 思维器官}{D. 效应器官}

\q 信息产业是指（\quad）
\choices{A. 只生产计算机硬件的产业}{B. 进行信息的收集、整理、存储、传输、处理及其应用服务的产业}
{C. 只提供通信线路和设备的产业}{D. 只从事软件开发和销售的产业}

\qtypename{二、多项选择题}{4 题，每题 4 分，共 16 分；每题有两个或两个以上正确项}
\q 关于数据与信息，下列说法正确的有（\quad）
\choicesline{A. 数据是信息载体}{B. 可消除不确定性}{C. 信息就是数据}{D. 数据是加工对象}

\q 信息的基本特点包括（\quad）
\choicesline{A. 可共享性}{B. 可存储性}{C. 可传输性}{D. 可消耗性}

\q 信息科学以（\quad）为理论基础。
\choicesline{A. 信息论}{B. 控制论}{C. 系统论}{D. 相对论}

\q 信息产业的特征包括（\quad）
\choicesline{A. 高渗透高催化}{B. 知识智力密集}{C. 高投入型}{D. 低增值型}

\qtypename{三、判断题}{6 题，每题 2 分，共 12 分；正确的画“√”，错误的画“×”并写出改正后的表述}
\q 数据是信息的载体，信息是数据经过加工后能为某个目的使用的内容。\quad（\quad）\quad 改正：\fixblank

\q 信息就是数据。\quad（\quad）\quad 改正：\fixblank

\q 系统论的奠基人是维纳，控制论的创立者是贝塔朗菲。\quad（\quad）\quad 改正：\fixblank

\q 广义信息论又称信息科学，研究范围比狭义信息论更广。\quad（\quad）\quad 改正：\fixblank

\q 科学是关于自然、社会和思维的发展与变化规律的知识体系，技术是运用科学规律实现某一目的的手段和方法。\quad（\quad）\quad 改正：\fixblank

\q 信息产业是低投入、低风险、低增值的产业。\quad（\quad）\quad 改正：\fixblank

\qtypename{四、配对题}{6 对，每对 1 分，共 6 分；把 A 栏与 B 栏对应内容配对}
\q 将下列概念与其对应内容配对，把 B 栏字母填入括号。
\par\vspace{0.2em}
\noindent\begin{tabular}{@{}p{0.44\textwidth}@{\hspace{1.2em}}p{0.44\textwidth}@{}}
（1）香农 \hfill（\quad） & A. 感觉器官 \\
（2）维纳 \hfill（\quad） & B. 思维器官 \\
（3）贝塔朗菲 \hfill（\quad） & C. 信息论创始人 \\
（4）传感技术 \hfill（\quad） & D. 控制论创始人 \\
（5）计算机技术 \hfill（\quad） & E. 系统论奠基人 \\
（6）三次信息 \hfill（\quad） & F. 综述、专题报告、年鉴 \\
\end{tabular}\par\vspace{0.3em}

\qtypename{五、简答题}{3 题，每题 8 分，共 24 分}
\q 简述数据、信息、知识与智慧的基本概念及相互关系，并结合农业生产举一例说明。\anslines{2.4cm}

\q 简述信息的基本特点（8 个方面）。\anslines{2.4cm}

\q 简述信息技术的四大范畴及其分别对应的人体器官，并说明信息科学的理论基础与主要工具。\anslines{2.4cm}

\newpage
\qtypename{六、综合题}{1 题，共 12 分}
\q 某县农业技术推广站把田间观测记录整理成数据库，又定期编写病虫害简报，最后形成年度专题报告。请判断上述三种材料分别属于一次、二次还是三次信息，并运用数据、信息、知识与智慧的关系，说明该站如何把原始记录逐步转化为可指导生产的决策支持。\anslines{5cm}

\end{document}
